\documentclass[10pt,a4paper,openany,twoside, eleve]{book}

\usepackage{layout_cours}
\setlist{leftmargin=5mm}

\begin{document}
\chapnumber{} % numéro du chapitre
\titre{Système linéaire à deux équations} % titre du chapitre
\classe{Seconde} % classe
\entetegal

% Debut du texte
% **************

\vspace{-5mm}
\section{Définition}

\begin{defi}{Système linéaire}{}
Résoudre le \textbf{système linéaire} $(S) : \syst{ax+by&=&c \\ a'x+b'y&=&c'}$ c'est trouver tous les couples de réels $(x;y)$ appelés \textbf{solutions du système} qui vérifient à la fois les deux équations.
\end{defi}

\begin{methode}{Vérifier une solution}{}
   Pour vérifier que le couple $(x;y)$ est une solution, je remplace les inconnues dans les deux équations et je calcule. On dit parfois que l'on \og réinjecte\fg dans l'équation.
\end{methode}

\begin{ex}{}{}
   Le couple $(2;1)$ est une solution de $(S_1) : \syst{2x+y&=&5 \\ x+4y&=&6}$ car :  \quad $\syst{2\times2+1&=&5 \\ 2+4\times1&=&6}$
\end{ex}

\begin{prop}{Nombre de solutions}{}
   Un système linéaire de deux équations à deux inconnues admet :
   \begin{itemize}
      \item Une unique solution $S=\{(x;y)\}$
      \item Aucune solution $S=\emptyset$, lorsque les deux équations sont \textbf{incompatibles}
      \item Une infinité de solutions, lorsque les deux équations sont \textbf{équivalentes}.
   \end{itemize}
\end{prop}

\begin{ex}{}{}
   \begin{itemize}
      \item Le système ci-dessus $(S_1) : \syst{2x+y&=&5 \\ x+4y&=&6}$ admet une unique solution $S=\left\{(2;1)\right\}$
      \item Le système $(S_2) : \syst{2x+y&=&5 \\ 2x+y&=&6}$ n'admet aucune solutions $S=\emptyset$ car les deux lignes sont incompatibles. Soit $2x+y=5$, soit $2x+y=6$, mais pas les deux en même temps !
      \item Le système $(S_3) : \syst{2x+y&=&5 \\ 6x-15&=&-3y}$ admet une infinité de solutions car les deux lignes sont équivalentes :
      
      $$6x-15=-3y \iff 6x+3y=15 \iff 2x+y=5$$
   \end{itemize}
\end{ex}

\vspace{-5mm}
\section{Méthodes de résolution par le calcul}
\vspace{-3mm}
\subsection{Résolution par substitution}

\begin{methode}{Résolution par substitution}{}
 Résoudre le système $(S):\syst{4x+y&=&7 \\ 3x-2y&=&8}$ par \textbf{substitution}. \\

\noindent Il s'agit d'exprimer une inconnue en fonction de l'autre à l'aide d'une des deux équations et de la remplacer par l'expression obtenue dans l'autre équation : \\

\begin{align*}
   (S) 
   &\Longleftrightarrow \syst{y&=&-4x+7 \\ 3x-2y&=&8} 
   &&\Longleftrightarrow \syst{y&=&-4x+7 \\ 3x-2(-4x+7)&=&8} 
   &&&\Longleftrightarrow \syst{y&=&-4x+7 \\ 3x+8x&=&8+14}\\
\end{align*}
\begin{align*}
   (S) 
   &\Longleftrightarrow \syst{y&=&-4x+7 \\ 11x&=&22} 
   &&\Longleftrightarrow \syst{y&=&-4\times 2+7 \\ x&=&2} 
   &&&\Longleftrightarrow \syst{y&=&-1 \\ x&=&2} \\
\end{align*}

\noindent D'où : $\mathcal{S}=\{(2;-1)\}$ \\
\end{methode}

\vspace{-4mm}
\begin{rem}{}{}
  La méthode par substitution n'est à utiliser (et encore \dots) que lorsqu'une inconnue s'exprime \emph{très facilement} en fonction de l'autre.
\end{rem}

\vspace{-3mm}
\subsection{Résolution par combinaison}

\begin{prop}{Opérations sur le système d'équation}{}
   Les opérations suivantes sur le système d'équation ne changent pas les solutions:
   \begin{itemize}
      \item Additionner deux lignes $(L_1) + (L_2)$
      \item Multiplier une ligne par un réel $a\in\R^*$ \textbf{non nul} $a\times(L_1)$
   \end{itemize}
\end{prop}

\begin{methode}{Résolution par combinaison linéaire}{}
Résoudre le système $(S):\syst{2x-3y&=&8\quad(L_1) \\ 5x+4y&=&-3\quad(L_2)}$ par \textbf{combinaison linéaire}. \\

\noindent On multiplie la première équation par $5$ et la deuxième par $(-2)$ de manière à obtenir le même coefficient devant l'inconnue $x$. Puis, on additionne les deux équations membre à membre : \\

$(S) \Longleftrightarrow \syst{10x-15y&=&40\quad(5\times L_1) \\ -10x-8y&=&6\quad(-2 \times L_2)} \Longleftrightarrow \syst{-23y&=&46 \quad (L_1+L_2) \\ -10x-8y&=&6}$\\

\noindent On obtient $y=-2$ avec la première équation et on remplace (méthode de substitution) dans la deuxième équation :\\ 

$(S) \Longleftrightarrow \syst{y&=&-2 \\ 2x-3\times(-2)&=&8} \Longleftrightarrow \syst{y&=&-2 \\ 2x&=&8-6} \Longleftrightarrow \syst{y&=&-2 \\ x&=&1}$  \\

\noindent D'où $\mathcal{S}=\{(1;-2)\}$ \\

\end{methode}


\vspace{-5mm}
\setcounter{section}{1}
\section*{Exercices}
\vspace{-3mm}
\begin{exercicesnofoot}
	\sectionexo{Nombre de solutions}
	\begin{exo}
      Quel est le nombre de solutions des systèmes suivants?

      \begin{multicols}{2}
         \begin{enumerate}
            \item $\syst{y &=& - 1,5 x + 2,4 \\ y &=& - 1,5 x - 8}$
            \item $\syst{4x -y = 8 \\ 2x = \frac{1}{2}y +4}$
            \item $\syst{y +5x &=& 2 \\ y -7x &=& 2}$
            \item $\syst{4x -y = 8 \\ 2x = -\frac{1}{2}y +4}$
         \end{enumerate}
      \end{multicols}
   \end{exo}
\vspace{-3mm}
\sectionexo{Résolution par substitution}
   \begin{exo}
      Résoudre les systèmes suivants par substitution et donner l'ensemble $S$ des couples $(x;y)$ de solutions.
      \begin{multicols}{2}
         \begin{enumerate}
            \item $\syst{x	-	2y	&=&	11 \\
2x	+	y	&=&	2}$
            \item $\syst{
               2x	-	y	&=&	6 \\
2y -x	&=&	3
            }$
            \item $\syst{
               x	+	2y	&=&	-2 \\
-x	-	3y	&=&	0
            }$
            \item $\syst{
               x-y&=&1 \\
               2x+3y&=&12
            }$
         \end{enumerate}
      \end{multicols}

   \end{exo}
	\sectionexo{Résolution par combinaison}
   \begin{exo}
      Résoudre les systèmes suivants par combinaison et donner l'ensemble $S$ des couples $(x;y)$ de solutions.
      \begin{multicols}{2}
         \begin{enumerate}
            \item $\syst{
               2x	+	3y	&=&	1 \\
x	-	3y	&=&	5
            }$
            \item $\syst{
4x	-	3y	&=&	11 \\
-x	+	2y	&=&	1 
            }$
            \item $\syst{
               -3x	-	2y	&=&	18 \\
               6x	+	5y	&=&	-42
            }$
            \item $\syst{
               2x+3y&=&4 \\
               x+5y&=&9
            }$
         \end{enumerate}
      \end{multicols}

   \end{exo}
\vspace{-3mm}
\sectionexo{Méthode libre}
   \begin{exo}
   
      Résoudre les systèmes suivants par la méthode de ton choix et donner l'ensemble $S$ des couples $(x;y)$ ou $(a;b)$ de solutions.
      \begin{multicols}{2}
      \begin{enumerate}
         \item $\syst{
y&=&4x-5 \\
y&=&-3x+2
         }$
         \item $\syst{
a-\frac{1}{3}b&=&3 \\
-3a+b&=&1
}$
         \item $\syst{
4y+3x&=&5 \\
6x+2y&=&1
         }$
         \item $\syst{
\frac{2}{3}a-4b&=&2 \\
-a+6b&=&-3
            }$
         \end{enumerate}
      \end{multicols}
   \end{exo}

\end{exercicesnofoot}

\newpage
\section*{Corrections}
\begin{multicols}{2}
	\raggedcolumns

\bfseries Exercice 2-1

\begin{align*}
   (S_1)
   &\Longleftrightarrow \syst{x - 2y&=&11 \\ 2x + y&=&2}  \\
   &\Longleftrightarrow \syst{x&=&11 + 2y \\ 2(11 + 2y) + y&=&2} \\
   &\Longleftrightarrow \syst{x&=&11 + 2y \\ 22 + 4y + y&=&2}  \\
   &\Longleftrightarrow \syst{x&=&11 + 2y \\ 5y&=&-20} \\
   &\Longleftrightarrow \syst{ x&=&3 \\ y&=&-4}
   \end{align*}

   $$S = \{ (3;-4) \}$$

\begin{align*}
   (S_2)
   &\Longleftrightarrow \syst{2x - y&=&6 \\ 2y - x&=&3} \\
   &\Longleftrightarrow \syst{y&=&2x - 6 \\ 2(2x - 6) - x&=&3} \\
   &\Longleftrightarrow \syst{y&=&2x - 6 \\ 4x - 12 - x&=&3} \\
   &\Longleftrightarrow \syst{y&=&2x - 6 \\ 3x&=&15} \\
   &\Longleftrightarrow \syst{y&=&4 \\ x&=&5 \\ }
   \end{align*}
   $$S = \{ (5;4) \}$$
      
\begin{align*}
   (S_3)
   &\Longleftrightarrow \syst{x + 2y&=&-2 \\ -x - 3y&=&0} \\
   &\Longleftrightarrow \syst{x&=&-2 - 2y \\ -(-2 - 2y) - 3y&=&0} \\
   &\Longleftrightarrow \syst{x&=&-2 - 2y \\ 2 + 2y - 3y&=&0} \\
   &\Longleftrightarrow \syst{x&=&-2 - 2y \\ 2 - y&=&0} \\
   &\Longleftrightarrow \syst{x&=&-6 \\ y&=&2}
   \end{align*}
   $$S = \{ (-6;2) \}$$
      
\begin{align*}
   (S_4)
   &\Longleftrightarrow \syst{x - y&=&1 \\ 2x + 3y&=&12} \\
   &\Longleftrightarrow \syst{x&=&1 + y \\ 2(1 + y) + 3y&=&12} \\
   &\Longleftrightarrow \syst{x&=&1 + y \\ 2 + 2y + 3y&=&12} \\
   &\Longleftrightarrow \syst{x&=&1 + y \\ 5y&=&10} \\
   &\Longleftrightarrow \syst{x&=&3 \\ y&=&2 }
   \end{align*}
   $$S = \{ (3;2) \}$$

\bfseries Exercice 3-1
\begin{align*}
   (S_1) 
   &\Longleftrightarrow \systl{2x + 3y &=& 1 &\quad (L_1) \\ x - 3y &=& 5 &\quad (L_2)} \\
   &\Longleftrightarrow \systl{2x + 3y &=& 1 &\quad (L_1) \\ 2x - 6y &=& 10 &\quad (L_2 \leftarrow 2 \times L_2)} \\
   &\Longleftrightarrow \systl{2x + 3y &=& 1 &\quad (L_1) \\ 0x - 9y &=& 9 &\quad (L_2 \leftarrow L_2 - L_1)} \\
   &\Longleftrightarrow \systl{2x +3\times(-1)&=& 1 &\quad (L_1) \\ y &=& -1 &\quad (L_2)} \\
   &\Longleftrightarrow \systl{x &=& 2 &\quad (L_1) \\ y &=& -1 &\quad (L_2)}
\end{align*}
   
   $$S = \left\{ \left(2\,;\,-1\right) \right\}$$

   \begin{align*}
      (S_2) 
      &\Longleftrightarrow \systl{4x - 3y &=& 11 &\quad (L_1) \\ -x + 2y &=& 1 &\quad (L_2)} \\
      &\Longleftrightarrow \systl{4x - 3y &=& 11 &\quad (L_1) \\ -4x + 8y &=& 4 &\quad (L_2 \leftarrow 4 \times L_2)} \\
      &\Longleftrightarrow \systl{4x - 3y &=& 11 &\quad (L_1) \\ 0x + 5y &=& 15 &\quad (L_2 \leftarrow L_1 + L_2)} \\
      &\Longleftrightarrow \systl{4x - 3\times\left(3\right) &=& 11 &\quad (L_1) \\ y &=& 3 &\quad (L_2 \leftarrow \frac{1}{5}L_2)} \\
      &\Longleftrightarrow \systl{x &=& 5 &\quad (L_1 \leftarrow \frac{1}{4}L_1) \\ y &=& 3 &\quad (L_2)}
      \end{align*}
      
      $$S = \left\{ (5\,;\,3) \right\}$$

      \begin{align*}
         (S_3) 
         &\Longleftrightarrow \systl{-3x - 2y &=& 18 &\quad (L_1) \\ 6x + 5y &=& -42 &\quad (L_2)} \\
         &\Longleftrightarrow \systl{-6x - 4y &=& 36 &\quad (L_1 \leftarrow 2 \times L_1) \\ 6x + 5y &=& -42 &\quad (L_2)} \\
         &\Longleftrightarrow \systl{-6x - 4y &=& 36 &\quad (L_1) \\ 0x + y &=& -6 &\quad (L_2 \leftarrow L_1 + L_2)} \\
         &\Longleftrightarrow \systl{-3x - 2\times(-6) &=& 18 &\quad (L_1) \\ y &=& -6 &\quad (L_2)} \\
         &\Longleftrightarrow \systl{-3x + 12 &=& 18 &\quad (L_1) \\ y &=& -6 &\quad (L_2)} \\
         &\Longleftrightarrow \systl{x &=& -2 &\quad (L_1 \leftarrow \frac{-1}{3}L_1) \\ y &=& -6 &\quad (L_2)}
         \end{align*}
         
         $$S = \left\{ (-2\,;\,-6) \right\}$$

         \begin{align*}
            (S_4) 
            &\Longleftrightarrow \systl{2x + 3y &=& 4 &\quad (L_1) \\ x + 5y &=& 9 &\quad (L_2)} \\
            &\Longleftrightarrow \systl{2x + 3y &=& 4 &\quad (L_1) \\ 2x + 10y &=& 18 &\quad (L_2 \leftarrow 2 \times L_2)} \\
            &\Longleftrightarrow \systl{2x + 3y &=& 4 &\quad (L_1) \\ 0x + 7y &=& 14 &\quad (L_2 \leftarrow L_2 - L_1)} \\
            &\Longleftrightarrow \systl{2x + 3\times(2) &=& 4 &\quad (L_1) \\ y &=& 2 &\quad (L_2 \leftarrow \frac{1}{7}L_2)} \\
            &\Longleftrightarrow \systl{2x + 6 &=& 4 &\quad (L_1) \\ y &=& 2 &\quad (L_2)} \\
            &\Longleftrightarrow \systl{x &=& -1 &\quad (L_1 \leftarrow \frac12 L_1) \\ y &=& 2 &\quad (L_2)}
            \end{align*}
            
            $$S = \left\{ (-1\,;\,2) \right\}$$
                     
      
\bfseries Exercice 4-1
\begin{align*}
   (S_1)
   &\Longleftrightarrow \syst{y&=&4x-5 \\ y&=&-3x+2} \\
   &\Longleftrightarrow \syst{4x-5&=&-3x+2 \\ y&=&4x-5} \\
   &\Longleftrightarrow \syst{7x&=&7 \\ y&=&4x-5} \\
   &\Longleftrightarrow \syst{x&=&1 \\ y&=&-1}
\end{align*}   

$$S=\left\{ (1\,;\,-1) \right\}$$

\columnbreak
\begin{align*}
   (S_2)
   &\Longleftrightarrow \systl{a-\frac{1}{3}b&=&3 &\quad (L_1)\\ -3a+b&=&1 &\quad (L_2)} \\
   &\Longleftrightarrow \systl{-3a+b&=&-9 &\quad (L_1\leftarrow -3L_1)\\ -3a+b&=&1 &\quad (L_2)} \\
   \end{align*}
   Les deux lignes sont incompatibles, le système n'a pas de solutions.

   $$S=\emptyset$$

   \begin{align*}
      (S_3)
      &\Longleftrightarrow \systl{4y+3x&=&5 &\quad (L_1)\\ 6x+2y&=&1 &\quad (L_2)}\\
      &\Longleftrightarrow \systl{6x+8y&=&10 &\quad (L_1\leftarrow 2L_1)\\  6x+2y&=&1 &\quad (L_2)}\\
      &\Longleftrightarrow \systl{0x+6y&=&9 &\quad (L_1\leftarrow L_1-L2)\\  6x+2y&=&1 &\quad (L_2)}\\
      &\Longleftrightarrow \systl{y&=&\frac{3}{2} &\quad (L_1\leftarrow \frac{1}{6}L_1)\\  6x+2\times(\frac{3}{2})&=&1 &\quad (L_2)}\\
      &\Longleftrightarrow \systl{y&=&\frac{3}{2} &\quad (L_1)\\  x&=&\frac{-1}{3} &\quad (L_2\leftarrow \frac{1}{6}L_2)}\\
      \end{align*}

$$S=\left\{ (\frac{-1}{3}\,;\,\frac{3}{2}) \right\}$$
      
      \begin{align*}
         (S_4)
         &\Longleftrightarrow \systl{\frac{2}{3}a - 4b&=&2 &\quad (L_1)\\ -a + 6b&=&-3 &\quad (L_2)} \\
         &\Longleftrightarrow \systl{-a + 6b&=&-3 &\quad (L_1 \leftarrow -\frac{3}{2}L_1)\\ -a + 6b&=&-3 &\quad (L_2)} \\
         &\Longleftrightarrow \begin{array}{rcl}  -a + 6b&=&-3   \end{array}
         \end{align*}
         
         Il y a donc une infinité de solutions, les deux lignes sont équivalentes.
         
	\end{multicols}
   \end{document}

